\section{Discussion}
\label{sec:discussion}

\paragraph{From colluding to adversarial.}
Definition~\ref{def:collusion} is behavioral: it flags a subset of agents that changed together. Whether the change is harmful is a separate question that requires a score. When a score $\sigma^{(t)}(A)$ on subsets of agents is available, a natural definition of an \emph{adversarial} subnetwork is one for which the score of the agents outside it, $\sigma^{(t)}(\N\setminus S)$, is lower than it would have been had $S$ kept its burn-in behavior; the difference is a treatment effect of $S$ on everyone else, and given $\hat S$ it can be estimated by extrapolating burn-in trends or by synthetic control from agents outside $\hat S$ in comparable systems. We have kept scores out of the detection problem deliberately, because on the platforms that motivate this work the score is observed only sometimes and is defined by the platform rather than by the observer, while the behavioral signal is always observed. A detected colluding subnetwork is the input to that assessment, and to any mitigation that follows; the interventions available to a platform operator (steering the environment or the score, quarantining edges, removing agents, resetting state) act on components the observer cannot touch, so in the vantablack setting the output of this work is a recommendation rather than an action.

\paragraph{Limitations.}
The burn-in assumption is load-bearing: a subnetwork that was already colluding during burn-in is absorbed into the baseline and is invisible to the test, as it is to adaptive control charts \citep{helm2026control}. We have treated the post-burn-in time $t$ as fixed. Estimating the onset $t^*$ is a sequential problem: a first-crossing procedure over $(t_0,t_\ell)$ pairs with a Bonferroni or permutation-calibrated family-wise level is the obvious construction, and it inherits the tradeoff of \citet{helm2026control}, since a fixed baseline is not evadable by slow defection but pays in delay while an adaptive baseline does the reverse. We leave onset estimation, and the characterization of delay against false-alarm rate, to future work. The joint permutation is exact under a null that also forbids a change in each agent's query distribution; when queries are self-selected, as on a platform, the conditional permutation of Section~\ref{sec:scan} is the appropriate null, at the cost of power when contexts are sparse. The candidate family is data-driven and of size $O(KN)$; a colluding subnetwork that is neither a cluster of the displacements nor a nested set along their leading direction, for instance two subnetworks moving in opposite directions, requires a richer family, and the scan's size standardization is what keeps a richer family from costing power. Finally, the embedding functions $g_\Q$ and $g_\X$ are fixed in advance; a colluding subnetwork whose shared change is invisible to the response embedding is invisible to the method, which is a limitation shared by every response-based representation.

\paragraph{Outlook.}
The method requires nothing from the agents beyond their observed behavior, and nothing from the platform beyond a log of who did what in which context. That is a low bar, and it is the bar that the incidents in Section~\ref{sec:intro} were observed to clear only in retrospect. We view detection of colluding subnetworks as one component of a monitoring stack for open multi-agent systems, alongside agent- and group-level change detection and system-level control charts, and the perspective-space representation as what lets these components share a single estimated object.
